%! Unit 1: Foundations --- Algebra 1 Review — Summative Assessment — Unit Test (blank)
\documentclass[10pt]{article}
\usepackage{algebra2-article}
\usepackage{algebra2-boxes}
\setlength{\workrowsep}{10pt}   % handwriting room; moves blank and key together

% Pre-drawn coordinate grid; #1 receives the plotted line and any overlay.
\newcommand{\gridplot}[1]{%
  \begin{tikzpicture}[x=0.30cm, y=0.19cm]
    \draw[step=1, linegray!45, very thin] (-5,-7) grid (5,7);
    \draw[->, charcoal, thick] (-5.4,0) -- (5.4,0) node[right, font=\scriptsize] {$x$};
    \draw[->, charcoal, thick] (0,-7.4) -- (0,7.4) node[above, font=\scriptsize] {$y$};
    \foreach \t in {-4,-2,2,4} {\node[below right=-2pt, font=\tiny] at (\t,0) {$\t$};}
    \foreach \t in {-6,-4,-2,2,4,6} {\node[above left=-2pt, font=\tiny] at (0,\t) {$\t$};}
    #1
  \end{tikzpicture}}

% Pre-drawn number line from -6 to 6; #1 receives the graphed solution.
\newcommand{\numline}[1]{%
  \begin{tikzpicture}[x=0.55cm, y=0.55cm]
    \draw[<->, charcoal, thick] (-6.7,0) -- (6.7,0);
    \foreach \t in {-6,-5,-4,-3,-2,-1,0,1,2,3,4,5,6} {
      \draw[charcoal] (\t,-0.16) -- (\t,0.16);
      \node[below=1pt, font=\tiny] at (\t,-0.16) {$\t$};}
    #1
  \end{tikzpicture}}

% Part-divider strip (headlinebox takes one arg: the background color)
\newcommand{\parthead}[1]{%
  \vspace{6pt}\noindent
  \begin{headlinebox}{forest}\color{white}\bfseries #1\end{headlinebox}%
  \vspace{4pt}}

\begin{document}

\pageheader{Unit 1: Foundations --- Algebra 1 Review}{Unit Test}
\namedateperiod

\vspace{0.06in}
\noindent\textbf{Instructions:} Show all work. No notes unless your teacher says so. A wrong answer
with a right idea can still earn credit; a right answer with no work may not. Each solve has ruled
space sized to the work it takes --- use it.

% ======================================================================
\boxguard[8]
\parthead{Part A --- Vocabulary (8 pts)}
Write the letter of the best description next to each term.

\vspace{4pt}
\noindent
\begin{minipage}[t]{0.34\textwidth}
\begin{enumerate}[leftmargin=2.2em, itemsep=7pt, label=\arabic*.]
  \item Rational number \blank{1.2cm}
  \item Associative property \blank{1.2cm}
  \item Factor completely \blank{1.2cm}
  \item Identity \blank{1.2cm}
  \item Zero product property \blank{1.2cm}
  \item Radicand \blank{1.2cm}
  \item Zero of a function \blank{1.2cm}
  \item Parallel lines \blank{1.2cm}
\end{enumerate}
\end{minipage}%
\hfill
\begin{minipage}[t]{0.63\textwidth}
\small
\begin{description}[leftmargin=1.4em, itemsep=3pt]
  \item[(A)] An equation true for every real number, so its solution set is all real numbers.
  \item[(B)] The expression written underneath the radical symbol.
  \item[(C)] Two distinct lines with equal slopes, so they never meet.
  \item[(D)] A real number that can be written as a ratio of two integers; its decimal terminates or repeats.
  \item[(E)] If a product equals $0$, then at least one factor equals $0$ --- the rule that turns a factored quadratic into solutions.
  \item[(F)] The property that lets you re-group terms without moving any of them past another.
  \item[(G)] To rewrite an expression as a product --- GCF first --- until no factor breaks down further over the integers.
  \item[(H)] An input whose output is $0$; also a solution of $f(x)=0$ and an $x$-intercept of the graph.
\end{description}
\end{minipage}

% ======================================================================
\boxguard[8]
\parthead{Part B --- Multiple Choice (12 pts)}
Circle the best answer.

\begin{enumerate}[leftmargin=1.9em, itemsep=9pt]
  \item Which property justifies the step $2\cdot(5x)=(2\cdot 5)x$?
    \\[3pt]\hspace*{1.2em} (A) Commutative property of multiplication
    \hspace{0.6cm} (B) Associative property of multiplication
    \\[2pt]\hspace*{1.2em} (C) Distributive property
    \hspace{3.15cm} (D) Multiplicative identity
  \item Simplified with no negative exponents, $\dfrac{28p^6q^{-1}}{7p^2q^3}$ is
    \hfill (A) $\dfrac{4p^4}{q^4}$ \quad (B) $\dfrac{4p^4}{q^2}$ \quad (C) $4p^4q^4$ \quad (D) $\dfrac{21p^4}{q^4}$
  \item The solution of $-5x+4\ge 24$ is
    \hfill (A) $x\ge -4$ \quad (B) $x\le -4$ \quad (C) $x\le 4$ \quad (D) $x\ge 4$
  \item How many \emph{real} solutions does $3x^2-2x+5=0$ have?
    \hfill (A) two \quad (B) one \quad (C) none \quad (D) infinitely many
  \item Which relation is \emph{not} a function?
    \\[3pt]\hspace*{1.2em} (A) $\{(-3,5),(0,5),(2,5),(6,5)\}$
    \hspace{0.9cm} (B) $\{(4,-2),(4,3),(7,1)\}$
    \\[2pt]\hspace*{1.2em} (C) $\{(-2,4),(-1,1),(0,0),(1,1)\}$
    \hspace{0.9cm} (D) $y=5x+2$
  \item \textbf{SOL-style.} Which equation represents the line shown?
    \begin{center}
    \begin{tabular}{@{}c l@{}}
    \gridplot{\draw[forest, very thick, domain=-5:5, smooth, variable=\x] plot ({\x},{0.5*\x+3});}
    &
    \begin{minipage}[c]{0.46\linewidth}
    \begin{enumerate}[label=(\Alph*), leftmargin=1.7em, itemsep=3pt]
      \item $y=2x+3$
      \item $y=\tfrac12 x+3$
      \item $y=-\tfrac12 x+3$
      \item $y=\tfrac12 x-3$
    \end{enumerate}
    \end{minipage}
    \\
    \end{tabular}
    \end{center}
\end{enumerate}

% ======================================================================
\boxguard[8]
\parthead{Part C --- Short Answer \& Computation (40 pts)}
Show your work.

\begin{enumerate}[leftmargin=1.9em, itemsep=10pt]
  \item \textbf{Numbers and justification.}
        \\[4pt]
        (a) Label each number \textbf{R} (rational) or \textbf{I} (irrational).
        \\[3pt]
        \hspace*{1.2em}$\sqrt{81}$ \blank{1.0cm} \quad $\sqrt{15}$ \blank{1.0cm} \quad
        $-\tfrac49$ \blank{1.0cm} \quad $0.\overline{6}$ \blank{1.0cm} \quad $\sqrt[3]{7}$ \blank{1.0cm}
        \\[4pt]
        (b) Name the property that justifies each step.
        \\[3pt]
        \hspace*{1.2em}$8(x+3)=8x+24$ \blank{4.6cm}
        \\[3pt]
        \hspace*{1.2em}$(2+x)+7=2+(x+7)$ \blank{4.6cm}
        \\[3pt]
        \hspace*{1.2em}$4x\cdot 1=4x$ \blank{4.6cm}
        \\[4pt]
        (c) ``The set of natural numbers is closed under division.'' Disprove it.
        \\[2pt]
        \writelines{2}

  \item \textbf{Translate and evaluate.}
        \\[4pt]
        (a) Write an expression for ``four less than the square of a number $n$.'' \blank{3.2cm}
        \\[4pt]
        (b) Evaluate $\dfrac{|a-b|}{4}+\sqrt{c}$ \; for $a=-6$, $b=6$, $c=36$.
        \begin{work}
          \frac{|a-b|}{4}+\sqrt{c} &= \frac{|-6-6|}{4}+\sqrt{36} \\
                                   &= \frac{|-12|}{4}+6 \\
                                   &= 3+6 \\
                                   &= 9
        \end{work}
        (c) Evaluate $\sqrt[3]{x}-6y$ \; for $x=-8$, $y=\tfrac13$.
        \begin{work}
          \sqrt[3]{x}-6y &= \sqrt[3]{-8}-6\left(\tfrac13\right) \\
                         &= -2-2 \\
                         &= -4
        \end{work}

  \item \textbf{Exponents and polynomial operations.}
        \\[4pt]
        (a) $(5m^4n)(-4m^2n^6)$
        \begin{work}
          (5m^4n)(-4m^2n^6) &= (5)(-4)\,m^{4+2}n^{1+6} \\
                            &= -20m^{6}n^{7}
        \end{work}
        (b) $\dfrac{42u^4v^2}{7u^4v^6}$ \; (no negative exponents)
        \begin{work}
          \frac{42u^4v^2}{7u^4v^6} &= 6u^{4-4}v^{2-6} \\
                                   &= 6u^{0}v^{-4} \\
                                   &= \frac{6}{v^{4}}
        \end{work}
        (c) $(4x-3)(2x+5)$
        \begin{work}
          (4x-3)(2x+5) &= 8x^2+20x-6x-15 \\
                       &= 8x^2+14x-15
        \end{work}
        (d) $(6x^2-2x+3)-(x^2+5x-4)$
        \begin{work}
          (6x^2-2x+3)-(x^2+5x-4) &= 6x^2-2x+3-x^2-5x+4 \\
                                 &= 5x^2-7x+7
        \end{work}

  \item \textbf{Factor completely.} GCF first, every time.
        \\[4pt]
        (a) $2x^2-50$
        \begin{work}
          2x^2-50 &= 2(x^2-25) \\
                  &= 2(x-5)(x+5)
        \end{work}
        (b) $x^2-4x-21$
        \begin{work}
          x^2-4x-21 &= (x-7)(x+3)
        \end{work}
        (c) $3x^2+11x+6$
        \begin{work}
          3x^2+11x+6 &= 3x^2+9x+2x+6 \\
                     &= 3x(x+3)+2(x+3) \\
                     &= (3x+2)(x+3)
        \end{work}
        (d) $5x^3-25x^2+30x$
        \begin{work}
          5x^3-25x^2+30x &= 5x(x^2-5x+6) \\
                         &= 5x(x-2)(x-3)
        \end{work}

  \item \textbf{Radicals and rational exponents.}
        \\[4pt]
        (a) $\sqrt{98}=$ \blank{2.4cm} \hspace{1.2cm}
        (b) $\sqrt[3]{24}=$ \blank{2.4cm}
        \\[4pt]
        (c) Write $\sqrt{x^5}$ with a rational exponent. \blank{2.6cm}
        \\[4pt]
        (d) Write $27^{2/3}$ as a radical, then evaluate it. \blank{4.6cm}
        \\[4pt]
        (e) $\sqrt{8}+\sqrt{32}=$
        \begin{work}
          \sqrt{8}+\sqrt{32} &= 2\sqrt{2}+4\sqrt{2} \\
                             &= 6\sqrt{2}
        \end{work}

  \item \textbf{Linear equations and inequalities.}
        \\[4pt]
        (a) Solve $5(2x-3)=6x+9$.
        \begin{work}
          5(2x-3) &= 6x+9 \\
           10x-15 &= 6x+9 \\
            4x-15 &= 9 \\
               4x &= 24 \\
                x &= 6
        \end{work}
        (b) Classify each solution set as \textbf{one solution}, \textbf{no solution}, or
        \textbf{infinitely many}. You should not need to finish solving.
        \\[3pt]
        \hspace*{1.2em}$2x-5=2x+1$ \blank{3.6cm} \hspace{0.8cm}
        $3(x+4)=3x+12$ \blank{3.6cm}
        \\[4pt]
        (c) Solve $V=\tfrac13 Bh$ for $h$.
        \begin{work}
             V &= \tfrac13 Bh \\
            3V &= Bh \\
          \frac{3V}{B} &= h
        \end{work}
        (d) Solve $-4x-3<9$, then graph the solution on the number line.
        \begin{work}
          -4x-3 &< 9 \\
            -4x &< 12 \\
              x &> -3
        \end{work}
        \begin{center}\numline{}\end{center}

  \item \textbf{Quadratic equations.}
        \\[4pt]
        (a) Solve by factoring: $x^2+2x-8=0$
        \begin{work}
            x^2+2x-8 &= 0 \\
          (x+4)(x-2) &= 0 \\
            x=-4 \quad&\text{or}\quad x=2
        \end{work}
        (b) Solve by square roots: $(x+2)^2=36$
        \begin{work}
          (x+2)^2 &= 36 \\
              x+2 &= \pm 6 \\
              x=4 \quad&\text{or}\quad x=-8
        \end{work}
        (c) Solve with the quadratic formula: $x^2+6x-3=0$ \; (exact answer)
        \begin{work}
          x &= \frac{-6\pm\sqrt{6^2-4(1)(-3)}}{2(1)} \\
            &= \frac{-6\pm\sqrt{36+12}}{2} \\
            &= \frac{-6\pm\sqrt{48}}{2} \\
            &= \frac{-6\pm 4\sqrt{3}}{2} \\
            &= -3\pm 2\sqrt{3}
        \end{work}
        (d) \emph{Without solving}, use the discriminant to say how many real solutions
        $x^2-10x+25=0$ has.
        \\[2pt] \hspace*{1.2em}Discriminant: \blank{3.4cm} \quad
        Number of real solutions: \blank{4.0cm}

  \item \textbf{Read the graph.} The linear function $g$ is graphed below; each gridline is one unit.
        \begin{center}
        \gridplot{\draw[forest, very thick, domain=-2.4:3.4, smooth, variable=\x] plot ({\x},{-2*\x+2});}
        \end{center}
        (a) $g(-2)=$ \blank{1.8cm} \hspace{1.2cm}
        (b) Solve $g(x)=0$: \; $x=$ \blank{1.8cm}
        \\[4pt]
        (c) That number in (b) is called the \blank{2.6cm} of $g$, and on the graph it is the
        point \blank{2.2cm}.
        \\[4pt]
        (d) Slope: \blank{1.8cm}; \quad $y$-intercept: \blank{2.0cm}; \quad
        equation of $g$: \blank{3.2cm}
        \\[4pt]
        (e) Is $g$ increasing or decreasing? \blank{2.6cm} \quad Domain: \blank{2.4cm}
        \\[4pt]
        (f) Write the equation of the line through $(0,-4)$ \textbf{parallel} to $g$
        \blank{3.2cm}, and the line through $(0,3)$ \textbf{perpendicular} to $g$
        \blank{3.2cm}.
\end{enumerate}

% ======================================================================
\boxguard[8]
\parthead{Part D --- Extended Response (12 pts)}
Explain your reasoning in full sentences.

\begin{enumerate}[leftmargin=1.9em, itemsep=14pt]
  \item \textbf{Three families, five questions.} Each table below is generated by a linear, a
        quadratic, or an exponential rule.
        \begin{center}
        \renewcommand{\arraystretch}{1.25}
        \begin{tabular}{@{}c@{\hspace{1.1cm}}c@{\hspace{1.1cm}}c@{}}
        \begin{tabular}{@{}c|c@{}}
          $x$ & $A(x)$ \\ \hline
          $0$ & $2$ \\ $1$ & $7$ \\ $2$ & $12$ \\ $3$ & $17$ \\
        \end{tabular}
        &
        \begin{tabular}{@{}c|c@{}}
          $x$ & $B(x)$ \\ \hline
          $0$ & $1$ \\ $1$ & $3$ \\ $2$ & $9$ \\ $3$ & $27$ \\
        \end{tabular}
        &
        \begin{tabular}{@{}c|c@{}}
          $x$ & $C(x)$ \\ \hline
          $0$ & $4$ \\ $1$ & $5$ \\ $2$ & $8$ \\ $3$ & $13$ \\
        \end{tabular}
        \\
        \end{tabular}
        \end{center}
        \textbf{(a)} Name the family of each table and give the numerical evidence --- constant
        difference, constant \emph{second} difference, or constant ratio.
        \\[2pt]
        \writelines{3}
        \textbf{(b)} Write a rule for each function.
        \\[4pt]
        \hspace*{1.2em}$A(x)=$ \blank{3.0cm} \quad $B(x)=$ \blank{3.0cm} \quad $C(x)=$ \blank{3.0cm}
        \\[6pt]
        \textbf{(c)} Which function is largest at $x=1$? At $x=4$? Show the values.
        \begin{work}
          A(1)=7, \quad B(1)&=3, \quad C(1)=5 \\
          A(4)=22, \quad B(4)&=81, \quad C(4)=20
        \end{work}
        \textbf{(d)} A classmate says ``$B$ is just the biggest one.'' Use your answers to (c) to
        explain why ``which is bigger'' is not a fair question until someone says \emph{where}.
        \\[2pt]
        \writelines{3}
        \textbf{(e)} If $B$ modeled a population of bacteria in hours, what domain would make sense,
        and why is that narrower than the domain of the rule by itself?
        \\[2pt]
        \writelines{3}

  \item \textbf{Modeling and solution sets.} Two shops quote a price for printing team shirts.
        \\[3pt]
        \hspace*{1.2em}\textbf{Ridgeway:} a one-time \$30 setup fee, then \$8 per shirt.
        \\[2pt]
        \hspace*{1.2em}\textbf{Lakeside:} no setup fee, \$14 per shirt.
        \\[4pt]
        \textbf{(a)} Write a function for the total cost at each shop for $s$ shirts.
        \\[4pt]
        \hspace*{1.2em}$R(s)=$ \blank{3.6cm} \hspace{1.4cm} $L(s)=$ \blank{3.6cm}
        \\[6pt]
        \textbf{(b)} Find the number of shirts at which the two costs are equal. Show every step.
        \begin{work}
          8s+30 &= 14s \\
             30 &= 6s \\
              5 &= s
        \end{work}
        Name the property of equality that justifies each step above.
        \\[3pt]
        \hspace*{1.2em}Step 1: \blank{5.4cm} \quad Step 2: \blank{5.4cm}
        \\[6pt]
        \textbf{(c)} Verify your answer by substituting it into \emph{both} functions.
        \begin{work}
          R(5) &= 8(5)+30 = 70 \\
          L(5) &= 14(5) = 70
        \end{work}
        \textbf{(d)} Which shop is cheaper for $3$ shirts? For $20$ shirts? Give the costs.
        \begin{work}
          R(3)=54, \quad L(3)&=42 \\
          R(20)=190, \quad L(20)&=280
        \end{work}
        \textbf{(e)} Suppose Lakeside changed its quote to a \$30 setup fee plus \$8 per shirt.
        Setting the two costs equal would give $8s+30=8s+30$. What is the solution set, what is
        that kind of equation called, and what does it mean about the two shops?
        \\[2pt]
        \writelines{3}
        \textbf{(f)} What domain makes sense for $s$ in this situation, and why?
        \\[2pt]
        \writelines{2}
\end{enumerate}

\end{document}
